\chapter{Emerging-Market FX and NDFs}\label{ch:m2:emerging-market-fx-and-ndfs}

For three years and four months the euro never cost less than 1.20 Swiss
francs. On 6 September 2011 the Swiss National Bank had declared that it would
``no longer tolerate'' a lower rate and was ``prepared to buy foreign currency in
unlimited quantities'', and the European Central Bank's daily reference rate,
1.1111 the day before, never again fell below 1.20. On 15 January
2015 the Swiss National Bank discontinued the floor. By the European Central
Bank's reference time that afternoon the euro was worth 1.0280 francs, 14.4\%
less than the day before, and a week later 0.9816. An online currency broker in
New York found that its customers owed it more than their deposits and borrowed
USD~300 million within a day to stay in business. Exchange rates that a
government fixes, bands, or controls behave calmly for years and then do not.
This chapter is about the currencies of emerging markets, and of a few rich
ones, where the state stands between the buyer and the seller: \omterm{def:m2:emerging-market-fx-and-ndfs:controls}{convertibility}
and \omterm{def:m2:emerging-market-fx-and-ndfs:controls}{capital controls}, the onshore and \omterm{def:m2:emerging-market-fx-and-ndfs:onshore}{offshore markets} they create, the
\omterm{def:m2:emerging-market-fx-and-ndfs:ndf}{non-deliverable forward} that the \omterm{def:m2:emerging-market-fx-and-ndfs:onshore}{offshore market} trades, and pegs, bands and the
day they break.

\section{Convertibility and capital controls}

\begin{definition}[Convertibility, capital control]\label{def:m2:emerging-market-fx-and-ndfs:controls}
A currency is \emph{convertible}\index{convertibility} when anyone may exchange
it for foreign currency, for any purpose, at a market rate. A \emph{capital
control}\index{capital control} is a legal restriction on such exchanges or on
the cross-border movement of money: limits on who may buy foreign currency, for
what, how much, or how quickly it may leave or enter the country.
\end{definition}

Most rich-country currencies are fully \omterm{def:m2:emerging-market-fx-and-ndfs:controls}{convertible}. Many emerging-market
currencies are \omterm{def:m2:emerging-market-fx-and-ndfs:controls}{convertible} for trade, a company paying for imports can buy
dollars, but not freely for investment: a foreign investor may need to register
its purchases of local bonds, and a local resident may be limited in how much it
can hold abroad. The controls exist to keep the exchange rate and domestic
credit under the authorities' control, and they have a direct consequence for
markets: the currency cannot be delivered outside the country, so foreigners who
want to hedge or speculate need an instrument that settles in something else.

\section{Onshore and offshore markets}

\begin{definition}[Onshore market, offshore market]\label{def:m2:emerging-market-fx-and-ndfs:onshore}
The \emph{onshore market}\index{onshore market} of a currency is the market in
its own country, under its regulators and controls; the \emph{offshore
market}\index{offshore market} is the market outside it, among non-residents,
where the currency trades on different terms or only in cash-settled form.
\end{definition}

The renminbi shows both at once. Onshore, the renminbi (code CNY) trades in the
interbank market around a central parity that the China Foreign Exchange Trade
System publishes each morning for the People's Bank of China; since 17 March
2014 the spot rate against the dollar may move up to 2\% above or below it. Hong
Kong, the pioneer offshore centre, trades the same currency outside those rules,
as CNH, with its own spot and forward prices and its own fixing, published by
the Treasury Markets Association. The two rates usually stay close, since
the flows between the markets are allowed in controlled channels, and diverge
when the controls bind (\cref{fig:m2:emerging-market-fx-and-ndfs:markets}).
Their forwards imply different renminbi interest rates: an onshore--offshore
basis.

\begin{omfigure}
\centering
\begin{tikzpicture}[font=\small,
  box/.style={draw,thick,rounded corners=2pt,align=center,minimum height=1.2cm,text width=4.6cm,inner sep=4pt},
  arr/.style={{Stealth[length=2.5mm]}-{Stealth[length=2.5mm]},thick}]
\node[box,draw=omThm,fill=omThm!8] (on) at (0,0) {\textbf{onshore (CNY)}\\interbank market, residents;\\spot within 2\% of the daily\\central parity};
\node[box,draw=omDef,fill=omDef!8] (off) at (7.2,0) {\textbf{offshore (CNH)}\\Hong Kong and other centres;\\deliverable renminbi,\\own spot, forwards and fixing};
\node[box,draw=omProp,fill=omProp!8] (ndf) at (7.2,-2.6) {\textbf{NDFs}\\cash-settled in dollars\\against the onshore fixing};
\draw[arr,dashed] (on) -- node[above,font=\footnotesize,align=center] {controlled\\channels} (off);
\draw[omProp,thick,-{Stealth[length=2.5mm]}] (on.south) |- node[pos=0.75,below,font=\footnotesize] {fixing (CNY SAEC)} (ndf.west);
\end{tikzpicture}

\omcaption{One currency, three markets. The onshore renminbi trades within a band
around the daily central parity; the offshore renminbi is deliverable outside
mainland China at its own prices; NDFs settle in dollars against the onshore
fixing. Money moves between the first two only through controlled channels, so
their prices can differ. Schematic.}\label{fig:m2:emerging-market-fx-and-ndfs:markets}
\end{omfigure}

\begin{example}[Two forward curves for one currency]\label{ex:m2:emerging-market-fx-and-ndfs:basis}
Suppose onshore USDCNY trades at 7.1000 spot and 7.0400 three months forward,
and offshore USDCNH at 7.1200 and 7.0700, with a dollar rate of 3.68\%. The
onshore forward implies a renminbi rate of 0.31\% for 91 days, the offshore one
0.89\%: a basis of 0.58 percentage points, the price of moving renminbi across the
line. The onshore band that day would run from 6.9580 to 7.2420.
\end{example}

\section{Non-deliverable forwards}

\begin{definition}[Non-deliverable forward, fixing source]\label{def:m2:emerging-market-fx-and-ndfs:ndf}
A \emph{non-deliverable forward}\index{non-deliverable forward} (NDF) is a
forward on a currency pair that is settled in cash, in dollars, for the
difference between the contract rate $K$ and a published fixing rate $S$ on a
fixing date: on a dollar notional $N$, the dollar buyer receives
$N(S - K)/S$. The \emph{fixing source}\index{fixing source} is the rate the
contract names for $S$, usually a rate published by the local central bank or
an official agent.
\end{definition}

The NDF lets a foreign investor hedge, or a fund speculate on, a currency it can
neither hold nor deliver: no local currency changes hands, only dollars, outside
the country. Its price follows the expected fixing and the dollar rate, and so
implies an offshore interest rate that can differ from the onshore one by the
basis of \cref{ex:m2:emerging-market-fx-and-ndfs:basis}. The industry's
standard terms name the \omterm{def:m2:emerging-market-fx-and-ndfs:ndf}{fixing source} for each currency and fix two business
days before settlement, as in the table below.

\begin{center}
\small
\begin{tabular}{@{}lll@{}}
\toprule
Currency & \omterm{def:m2:emerging-market-fx-and-ndfs:ndf}{Fixing source} (EMTA code) & Published by, around \\
\midrule
Indian rupee & INR FBIL & Financial Benchmarks India, 13:30 Mumbai \\
Brazilian real & BRL PTAX & Banco Central do Brasil, 13:15 S\~ao Paulo \\
Chinese renminbi & CNY SAEC & CFETS for the People's Bank, 9:15 Beijing \\
Korean won & KRW KFTC18 & Seoul Money Brokerage, 16:00 Seoul \\
Taiwan dollar & TWD TAIFX1 & Taipei Forex, 11:00 Taipei \\
\bottomrule
\end{tabular}
\par\smallskip
{\footnotesize \omterm{def:m2:emerging-market-fx-and-ndfs:ndf}{Fixing sources} of five NDF currencies, each ``for settlement in two
business days'', in Annex A to the 1998 FX and Currency Option Definitions (as of
June 2023).}
\end{center}

\begin{omfigure}
\centering
\begin{tikzpicture}[font=\small,x=1.35cm]
\draw[thick,-{Stealth}] (0,0) -- (9.6,0) node[right,font=\footnotesize] {time};
\foreach \x/\l in {0/trade: agree $K$,6.4/fixing date: $S$ published,8.4/settlement} {\draw (\x,0.12) -- (\x,-0.12); \node[below,font=\footnotesize,align=center] at (\x,-0.15) {\l};}
\draw[gray,decorate,decoration={brace,amplitude=4pt}] (6.4,0.35) -- (8.4,0.35) node[midway,above=4pt,font=\footnotesize,gray] {two business days};
\draw[omThm,thick,-{Stealth}] (8.4,-1.2) -- (8.4,-2.0);
\node[omThm,font=\footnotesize,align=center] at (8.4,-2.45) {one payment in dollars:\\$N(S-K)/S$ to the dollar buyer};
\node[omDef,font=\footnotesize,align=center] at (3.2,-1.6) {no local currency changes hands;\\nothing needs to enter or leave the country};
\end{tikzpicture}

\omcaption{The life of a \omterm{def:m2:emerging-market-fx-and-ndfs:ndf}{non-deliverable forward}. The contract rate is agreed at
the trade; two business days before settlement the \omterm{def:m2:emerging-market-fx-and-ndfs:ndf}{fixing source} publishes the
rate, and on the settlement date the difference is paid in dollars.
Schematic.}\label{fig:m2:emerging-market-fx-and-ndfs:ndf}
\end{omfigure}

\begin{example}[A won NDF]\label{ex:m2:emerging-market-fx-and-ndfs:krw}
A fund buys USD~10 million three months forward against won in an NDF at 1\,380.00,
settling on Tuesday 13 October 2026. The fixing is taken two business days
before, on Friday 9 October, at 1\,425.50: the fund receives
$10\,000\,000 \times 45.50/1\,425.50 = \text{USD}~319\,186$. Had the won fixed at
1\,350.00, the fund would have paid USD~222\,222.
\end{example}

\section{Pegs, bands and interventions}

\begin{definition}[Currency peg, currency band, FX intervention]\label{def:m2:emerging-market-fx-and-ndfs:peg}
A \emph{currency peg}\index{currency peg} is a policy of holding the exchange
rate at a fixed level against another currency; a \emph{currency
band}\index{currency band} lets it move within set limits around a central
rate. An \emph{FX intervention}\index{FX intervention} is the purchase or sale
of foreign currency by the authorities to move or hold the rate.
\end{definition}

Hong Kong's dollar has been linked to the US dollar since 17 October 1983; its
monetary authority keeps it within a band of 7.75 to 7.85 per dollar through a
currency board, which stands ready to exchange the two currencies at the edges of
the band. A band like the renminbi's is managed by the daily
central parity and by intervention when the rate approaches its limits. A
one-sided floor like the Swiss franc's commits the central bank to sell its own
currency without limit, which it can always do, since it can create it; a
ceiling commits it to sell foreign currency, which it can do only as long as its
reserves last.

\begin{dated}{2026-09}{Pegs, bands and fixings}\label{dat:m2:emerging-market-fx-and-ndfs:regimes}
Hong Kong dollar: HK\$7.75 to 7.85 per US dollar under the Linked Exchange Rate
System, in place since 17 October 1983. Renminbi: onshore spot within 2\% of the
central parity, the width set from 17 March 2014; NDF fixing CNY SAEC at about
9:15 Beijing time, offshore fixing CNY CNHHK at about 11:30 Hong Kong time
(standard terms of June 2023). Swiss franc: no minimum rate since 15 January
2015.
\end{dated}

For a trader, a credible peg makes the rate nearly riskless and its options
nearly worthless, and that is the trap: implied volatility at the peg measures
the market's confidence in the policy, not the size of the move if the policy
ends. The move, when it comes, is a jump.

\section{The franc unpeg}

\begin{omfigure}
\begin{tikzpicture}
\begin{axis}[width=11cm,height=5cm,
  ylabel={Swiss francs per euro},
  tick label style={font=\small},label style={font=\small},
  xmin=2010,xmax=2017,ymin=0.95,ymax=1.5,grid=major,grid style={gray!25},
  xtick={2010,2011,2012,2013,2014,2015,2016,2017},xticklabel style={/pgf/number format/1000 sep={}},
  table/col sep=comma]
\addplot[omThm,thick] table[x=t,y=eurchf]{figdata/markets-2/18-emerging-market-fx-and-ndfs/eurchf.csv};
\draw[omProp,dashed] (axis cs:2011.68,1.2) -- (axis cs:2015.04,1.2);
\node[omProp,font=\footnotesize,anchor=north,align=center] at (axis cs:2013.3,1.185) {floor at 1.20\\6 Sep 2011 to 15 Jan 2015};
\end{axis}
\end{tikzpicture}

\omcaption{The euro in Swiss francs, ECB daily reference rates, 2010 to 2016.
After the franc's rise of 2010--2011 the Swiss National Bank held the rate at or
above 1.20 for over three years; on 15 January 2015 it let it go, and the
reference rate fell from 1.2010 to 1.0280 in a day. Data: ECB Data Portal,
series EXR.D.CHF.EUR.SP00.A.}\label{fig:m2:emerging-market-fx-and-ndfs:eurchf}
\end{omfigure}

The Swiss floor was one-sided: the Swiss National Bank promised to sell francs
at 1.20 without limit, and for three years it bought euros whenever the market
pushed the rate down. A trader could treat the floor as a wall: buying euros
near 1.20, with a stop order a little below, looked like a trade with little to
lose, and could be held with leverage. When the central bank withdrew, on 15
January 2015, and cut its deposit rate to $-0.75\%$ at the same time, the rate
did not pass through the levels in between; it jumped. A stop order becomes a
market order when its level is traded through and is filled at the next price
available, which that day could be far below it. Leveraged clients could lose
more than their deposits, and their brokers the rest.

\begin{proposition}[Leverage and a jump]\label{prop:m2:emerging-market-fx-and-ndfs:leverage}
A position of $L$ times the margin deposited loses $L\,m$ times the margin when
the price jumps by $m$ against it before any stop can be executed; the client's
balance becomes negative when $m > 1/L$, and the excess is owed to the broker.
\end{proposition}

\begin{proof}
With margin $M$ and position $LM$, a relative move $m$ loses $LMm$, which exceeds
$M$ exactly when $Lm > 1$.
\end{proof}

A 14.4\% jump turns 20-to-1 leverage into a loss of 2.88 times the margin and
50-to-1 into 7.2 times (\cref{fig:m2:emerging-market-fx-and-ndfs:leverage}). The
broker of the hook, FXCM, reported on 19 January 2015 that its clients' debit
balances had obliged its regulated subsidiaries to raise capital, and that it
had borrowed USD~300 million from Leucadia on 16 January for two years at an
initial 10\% a year, rising to as much as 17\%, to cover them.

\begin{omfigure}
\begin{tikzpicture}
\begin{axis}[width=11cm,height=4.4cm,ybar,bar width=22pt,
  xlabel={leverage (position / margin)},ylabel={loss / margin},
  tick label style={font=\small},label style={font=\small},
  xmin=-0.6,xmax=3.6,ymin=0,ymax=9,grid=major,grid style={gray!25},
  xtick={0,1,2,3},xticklabels={5,10,20,50},
  nodes near coords,every node near coord/.append style={font=\scriptsize},
  table/col sep=comma]
\addplot[fill=omProp!60,draw=omProp] table[x=k,y=multiple]{figdata/markets-2/18-emerging-market-fx-and-ndfs/leverage.csv};
\draw[omThm,thick,dashed] (axis cs:-0.6,1) -- (axis cs:3.6,1);
\node[omThm,font=\scriptsize,anchor=south] at (axis cs:0.5,1.05) {deposit lost};
\end{axis}
\end{tikzpicture}

\omcaption{Loss on a long EURCHF position, as a multiple of the margin deposited,
for the move of the ECB reference rate from 1.2010 on 14 January to 1.0280 on 15
January 2015, if no stop was filled before the afternoon price. Above the dashed
line the client owes the broker money. Data: ECB reference rates; the chapter's
tutorial.}\label{fig:m2:emerging-market-fx-and-ndfs:leverage}
\end{omfigure}

\section{Tutorial: NDF settlement and the onshore--offshore basis}\label{tut:m2:emerging-market-fx-and-ndfs:ndf}

\begin{tutorial}
\textbf{Goal.} Settle an NDF against its fixing, date the fixing, compute the
interest rates implied by onshore and offshore forwards and their basis, and the
loss of a leveraged position on a jump. \textbf{End state:}
\cref{ex:m2:emerging-market-fx-and-ndfs:basis,ex:m2:emerging-market-fx-and-ndfs:krw},
\cref{fig:m2:emerging-market-fx-and-ndfs:eurchf,fig:m2:emerging-market-fx-and-ndfs:leverage}
and the numbers of the weekend problem.
\begin{tutsteps}
\item \textbf{The NDF}: its settlement and its fixing date.
\omcode{code/firm/ndf/firm_ndf.py}{11}{24}{NDF settlement and fixing date.}\label{lst:m2:emerging-market-fx-and-ndfs:ndf}
\item \textbf{The basis}: implied local rates from each market's forward.
\omcode{code/firm/ndf/firm_ndf.py}{27}{36}{Implied local-currency rate and the onshore--offshore basis.}\label{lst:m2:emerging-market-fx-and-ndfs:basis}
\item \textbf{Run} \texttt{ndf\_demo.krw\_ndf()}, \texttt{ndf\_demo.cny\_cnh()},
\texttt{ndf\_demo.unpeg()} and \texttt{fig\_ndf.py}.
\end{tutsteps}
\textbf{What to change next.} Price the NDF's value before fixing, from the NDF
forward curve and dollar discounting; then compute how the basis of
\cref{ex:m2:emerging-market-fx-and-ndfs:basis} changes the value of a three-month
hedge for a foreign holder of Chinese bonds.
\end{tutorial}

\section{Build: the NDF calculator}\label{bld:m2:emerging-market-fx-and-ndfs:ndf}

\begin{build}
\bfield{Purpose} The miniature firm trades emerging-market currencies offshore:
it needs NDF settlement amounts and fixing dates, the offshore rates implied by
NDF curves, and a measure of what leverage costs on a jump.
\bfield{Interface} \texttt{ndf\_settlement(notional\_usd, contract, fixing,
buyer)}; \texttt{fixing\_date(settlement, holidays, lag)};
\texttt{implied\_local\_rate}; \texttt{onshore\_offshore\_basis};
\texttt{band(central, width)}; \texttt{leveraged\_loss(move, leverage)}.
\bfield{Rules} Rates as local currency per dollar; settlement in dollars;
fixing two business days before settlement by default; dollar rates on
actual/360 and local rates on actual/365 unless told otherwise.
\bfield{Acceptance tests} \texttt{code/firm/ndf/tests/}: zero settlement at the
contract rate and the signs either side; a fixing date that skips a weekend and a
holiday; an implied rate recovered from a forward built with it; the band's
limits; the leverage multiple.
\bfield{Stretch} Fallbacks when the \omterm{def:m2:emerging-market-fx-and-ndfs:ndf}{fixing source} is not published (the
disruption rules of the standard terms); NDF curves bootstrapped from quotes;
the value of an option on a peg's break, priced with a jump.
\end{build}

\begin{omsources}
\item Swiss National Bank, press releases of 6 September 2011 and 15 January 2015;
European Central Bank, euro reference rates (Data Portal).
\item FXCM Inc., Form 8-K, 19 January 2015 (Exhibit 99.1).
\item People's Bank of China, announcement on the renminbi trading band, March
2014; Hong Kong Monetary Authority, Linked Exchange Rate System, and renminbi
business booklet.
\item EMTA and ISDA, Annex A to the 1998 FX and Currency Option Definitions, June
2023.
\end{omsources}

\section{Exercises}

\begin{exercise}[$\star$]\label{exo:m2:emerging-market-fx-and-ndfs:1}
A dollar seller in the won NDF of \cref{ex:m2:emerging-market-fx-and-ndfs:krw}:
what does it receive or pay at a fixing of 1\,425.50, and at 1\,350.00?
\end{exercise}

\begin{exercise}[$\star$]\label{exo:m2:emerging-market-fx-and-ndfs:2}
Why does a foreign fund hedge Indian rupee bonds with an NDF rather than a
deliverable forward?
\end{exercise}

\begin{exercise}[$\star$]\label{exo:m2:emerging-market-fx-and-ndfs:3}
The renminbi's central parity is 7.1000. Give the onshore trading band, and the
Hong Kong dollar's band.
\end{exercise}

\begin{exercise}[$\star\star$]\label{exo:m2:emerging-market-fx-and-ndfs:4}
An NDF settles on Monday 5 October 2026 and Friday 2 October is a holiday for its
\omterm{def:m2:emerging-market-fx-and-ndfs:ndf}{fixing source}. When is the fixing taken?
\end{exercise}

\begin{exercise}[$\star\star$]\label{exo:m2:emerging-market-fx-and-ndfs:5}
In \cref{ex:m2:emerging-market-fx-and-ndfs:basis}, what does a positive
offshore-minus-onshore basis mean for a foreign investor hedging renminbi
bonds offshore?
\end{exercise}

\begin{exercise}[$\star\star$]\label{exo:m2:emerging-market-fx-and-ndfs:6}
Why is a floor, as the Swiss franc's was, easier to defend than a ceiling?
\end{exercise}

\begin{exercise}[$\star\star\star$]\label{exo:m2:emerging-market-fx-and-ndfs:7}
\emph{Coding.} With the ECB data, give the lowest EURCHF reference rate between
7 September 2011 and 14 January 2015, the reference rates of 14 and 15 January
2015, and the lowest of January 2015.
\end{exercise}

\begin{exercise}[$\star\star\star$]\label{exo:m2:emerging-market-fx-and-ndfs:8}
\emph{Find the flaw.} ``EURCHF one-month implied volatility is 1\% with the
floor in place: a long EURCHF position with a stop 2\% lower risks at most 2\%.''
Correct it.
\end{exercise}

\section{Problem: The Floor that Broke}

\begin{problem}[{Weekend problem --- a leveraged client on 15 January 2015}]\label{pb:m2:emerging-market-fx-and-ndfs:1}
A retail client is long EUR~1 million against francs at 1.2010, the reference rate
of 14 January 2015, with 20-to-1 leverage and a stop at 1.19. On 15 January the
floor is removed; the reference rate that afternoon is 1.0280.

\medskip\noindent\textbf{Part I --- The position.}
\begin{enumerate}
\item Give the margin deposited, in francs.
\item What loss would the stop at 1.19 have produced?
\item Give the move to 1.0280 in percent.
\item Give the loss at 1.0280, in francs and as a multiple of the margin.
\item Give the client's balance.
\end{enumerate}

\medskip\noindent\textbf{Part II --- Why the stop failed.}
\begin{enumerate}[resume]
\item What is a stop order, and what does it guarantee?
\item Why were there no buyers of euros near 1.19?
\item Why did the floor's calm make the leverage look safe?
\item What did one-month implied volatility say before, and what did it miss?
\item How would a put option on EURCHF have behaved?
\end{enumerate}

\medskip\noindent\textbf{Part III --- The broker.}
\begin{enumerate}[resume]
\item Who bears the client's negative balance at first?
\item How did FXCM cover its clients' negative balances?
\item What did the loan cost?
\item How should a broker set leverage on a pegged pair?
\item What protects retail clients from negative balances in some jurisdictions?
\end{enumerate}

\medskip\noindent\textbf{Part IV --- Judgement.}
\begin{enumerate}[resume]
\item Why do central banks abandon pegs suddenly rather than gradually?
\item What other pegs or bands could break the same way?
\item How would you size a position in a pegged currency?
\item State the \emph{named result}: the loss on the 20-to-1 long at the ECB
reference rate of 15 January 2015.
\item In one sentence: what is the risk of a pegged currency?
\end{enumerate}
\end{problem}

\section{Interview questions}

\begin{interviewq}[$\star$ \iqroles{trader, developer}]\label{iq:m2:emerging-market-fx-and-ndfs:1}
What is an NDF, and how is it settled?
\end{interviewq}

\begin{interviewq}[$\star$ \iqroles{trader, researcher}]\label{iq:m2:emerging-market-fx-and-ndfs:2}
Why can the onshore and offshore renminbi trade at different prices?
\end{interviewq}

\begin{interviewq}[$\star\star$ \iqroles{researcher}]\label{iq:m2:emerging-market-fx-and-ndfs:3}
How would you price an option on a currency whose peg may break?
\end{interviewq}

\begin{interviewq}[$\star\star$ \iqroles{bank, trader}]\label{iq:m2:emerging-market-fx-and-ndfs:4}
What went wrong for leveraged traders and brokers on 15 January 2015?
\end{interviewq}

\begin{interviewq}[$\star\star$ \iqroles{trader}]\label{iq:m2:emerging-market-fx-and-ndfs:5}
A client asks to hedge USD~100 million of Brazilian bonds. How do you structure
it, and what can go wrong at the fixing?
\end{interviewq}

\begin{interviewq}[$\star\star\star$ \iqroles{developer, bank}]\label{iq:m2:emerging-market-fx-and-ndfs:6}
Design the risk system of a retail FX broker so that a move like that of 15
January 2015 does not threaten its solvency.
\end{interviewq}
